\documentclass[11pt]{tlc-article}
\begin{document}
% {{{ Document Title, Author, and Abstract
  \tlcTitlePageAndTableOfContents
  {Document Header and Footer Design}
  {\autodocDepartmentName}
  {This document describes the design used by the Automated Documentation
    Framework to produce consistent headers and footers for all documents. The
    header and footer reserve layout space for the \gls{company} logo,
    document title, document number, document date, and other fields required by
    the \gls{qms}. The implementation details describe the steps needed to
    satisfy the design objectives.}

% -------------------------------------------------------------------------- }}}
% {{{ File names

\section{File names}
This section defines the symbolic names used by the directory tree and matrices.
The file symbols identify the document customization files, and the superscript
symbols identify where each file is located. Defining these terms first keeps
the later transformation rules compact and repeatable while supporting
\gls{configuration-identification}.

\tlcVspace\

\begin{description}
  \item[\fa] = additional-layout.tex
  \item[\fb] = header-footer.tex
  \item[\fc] = logo.png
  \item[\fd] = version.tex
\end{description}

\begin{description}
  \item[\ra] = \raName
  \item[\rb] = \rbName
  \item[\rc] = \rcName
  \item[\rd] = \rdName
  \item[\re] = autoheaderfooter
\end{description}

% -------------------------------------------------------------------------- }}}
% {{{ Directory Structure

\section{Directory Structure}
This section defines the \gls{document-structure} used to resolve reusable
header and footer assets. Each report can provide local data files, reports
inside \texttt{project-core} can also inherit shared data files, and all reports
can fall back to autoheaderfooter defaults. Given the directory structure below,
transform matrix \ref{tab:src} to matrix \ref{tab:dst}.

\tlcVspace\

\dirtree{%
.1 repo-root.
.2 autoheaderfooter~\re.
.3 data.
.4 \faName~\fea.
.4 \fbName~\feb.
.4 \fcName~\fec.
.4 \fdName~\fed.
.2 project-report~\ra.
.3 data.
.4 \faName~\faa.
.4 \fbName~\fab.
.4 \fcName~\fac.
.4 \fdName~\fad.
.2 project-core.
.3 reports.
.4 shared.
.5 data.
.6 \faName~\fba.
.6 \fbName~\fbb.
.6 \fcName~\fbc.
.6 \fdName~\fbd.
.4 test-report.
.5 data.
.6 \faName~\fca.
.6 \fbName~\fcb.
.6 \fcName~\fcc.
.6 \fdName~\fcd.
.4 verification-report.
.5 data.
.6 \faName~\fda.
.6 \fbName~\fdb.
.6 \fcName~\fdc.
.6 \fdName~\fdd.
}

% -------------------------------------------------------------------------- }}}
% {{{ Source and Transformed Matrix

\section{Source Matrix}
This section maps the directory tree to a compact matrix form. Each row
represents a \gls{report} or shared/default data location, and each column
represents one file type. The symbols in this matrix are source inputs before
inheritance and fallback rules are applied.

\tlcVspace\

\begin{table}[ht]
  \begin{center}
    \begin{tabular}{l|c|c|c|c}
              & \fa  & \fb  & \fc  & \fd  \\ \hline
      \raName & \faa & \fab & \fac & \fad \\ \hline
      \rbName & \fba & \fbb & \fbc & \fbd \\ \hline
      \rcName & \fca & \fcb & \fcc & \fcd \\ \hline
      \rdName & \fda & \fdb & \fdc & \fdd \\ \hline
    \end{tabular}
    \caption{Source Matrix}\label{tab:src}
  \end{center}
\end{table}

\section{Transformed Matrix}
This section shows the matrix after the composition rules are applied. Each
\gls{report} output cell lists the files that participate in the resolved value
for that report and file type. The \rbName\ row is marked \nota\ because shared
files are inputs to report-specific compositions, not a standalone report
output.

\tlcVspace\

\begin{table}[ht]
  \begin{center}
    \begin{tabular}{l|c|c|c|c}
              & \fa   & \fb   & \fc   & \fd   \\ \hline
      \raName & \faaa & \faab & \faac & \faad \\ \hline
      \rbName & \nota & \notb & \notc & \notd \\ \hline
      \rcName & \fcca & \fccb & \fccc & \fccd \\ \hline
      \rdName & \fdda & \fddb & \fddc & \fddd \\ \hline
    \end{tabular}
    \caption{Transformed Matrix}\label{tab:dst}
  \end{center}
\end{table}

% -------------------------------------------------------------------------- }}}
% {{{ Transformation Explained

\section{Transformation Explained}

This section explains how the source matrix is converted into the transformed
matrix. The same rule is applied independently to each file type, so
\faName, \fbName, \fcName, and \fdName\ each compose with the matching shared
and default file when those files are applicable.

\tlcVspace\

Each report-specific data file is resolved by looking first in the report's
local data folder, then in the shared report data folder when the report belongs
to \texttt{project-core}, and finally in the autoheaderfooter defaults. This is
a \gls{configuration-management} rule for document layout inputs. The \raName\
\gls{report} is outside \texttt{project-core}, so its local files compose
directly with autoheaderfooter defaults. The \rbName\ row is not transformed
into a standalone output because shared files only participate in other reports'
compositions.

\tlcVspace\

Using \faName=\fdda~ as the example, one can trace the potential path from the
\rdName\ report file, to the shared report file, to the autoheaderfooter default
file. The \faName~ file specializes customizations that are shared by documents.
Similarly, \fdName~ represents the header and footer for each document. The
\gls{design-verification} goal is to define it once and reuse it across
reports.

\section{Example \rdName~ report using a composite \fdName}
\subsection{Default Version Information}

This example follows the \rdName\ report's \fdName\ cell in the transformed
matrix. The example is representative because it includes all three levels:
report-local data, shared report data, and the autoheaderfooter default.

\tlcVspace\

\faName~ has an initial implementation that is consistent with the default
layout provided by tlc-article. The header and footer use the following
placement rules to complete the document layout specification.

\begin{description}
  \item[lhead] The company logo is written to the left-side header.
  \item[chead] The document title is written to the center header.
  \item[rhead] The document status, date, and version is written to the right
    side of the header.
  \item[lfoot] The institution name is written to the left-side footer.
  \item[cfoot] Permission for use is written to the center footer.
  \item[rfoot] Page \# of Pages \# is written to the right-side footer.
\end{description}

\subsection{\fdName}
This subsection identifies the version and status fields that flow into
\fdName. The default implementation reuses attributes known to tlc-article, and
report-specific files override only the values that need to differ for a
particular \gls{report}.

\tlcVspace\

\begin{description}
  \item[hfVersion] \hfVersion
  \item[hfDate] \hfDate
  \item[hfStatus] \hfStatus
  \item[hfInstitution] \hfInstitution
  \item[hfPermission] \hfPermission
\end{description}

\subsection{Composition}
This subsection lists the include order for the \fddd\ composition. Reading the
chain from local to default shows which file has the final opportunity to
specialize each value and which file provides the baseline when no override is
present.

\tlcVspace\

Given \fddd:
\begin{description}
  \item[\fdd] includes \fbd~ first and overrides hfVersion.
  \item[\fbd] includes \fed~ first and may override hfDate and hfStatus.
  \item[\fed] defines defaults for hfVersion, hfDate, hfStatus,
    hfInstitution, and hfPermission.
\end{description}

% -------------------------------------------------------------------------- }}}
% {{{ Print acronyms and glossaries

\printGlossaries%

% -------------------------------------------------------------------------- }}}
% {{{ Change summary

\input{change-summary.tex}

% -------------------------------------------------------------------------- }}}
\end{document}
